% Discussion. Labels defined here use the prefix "discussion:".
\section{Discussion}\label{sec:discussion}

The results separate the size of a proof from the search and local work used
to obtain it. They identify several reasons why a specified relaxation can
require a large certificate, and several reasons why a search can be large
when a small certificate exists. This distinction helps interpret both the
mathematical examples and the archived computations.

\subsection{Interpreting a certificate obstruction}\label{discussion:obstruction}

For relaxations with a positive quadratic vertex gap, the covering profile
in \cref{sec:geometry} measures how many near-optimal points must be placed
near test-box vertices. With the hypotheses for the matching upper bound,
it determines spatial certificate size up to a logarithmic factor.
Quadratic growth around a regular positive-dimensional optimal set gives a
polynomial rate, while an isolated degenerate minimum can give a polynomial
rate for a different reason. The all-face characterization in
\cref{sec:face-exact} further shows that boundary geometry matters: a proper
root face can govern the exponent even when the full-dimensional singularity
has a smaller exponent.

These are obstructions for the stated relaxation. A choice of branching
rule cannot produce a certificate below its minimum size. A stronger
relaxation, a change of factorization, or a different certificate format can
change that minimum. Tightening must be assessed with the operation costs:
under the pointwise gap assumptions it preserves a lower bound on charged
certificate events, while a node lower bound additionally requires bounded
rounds and appropriate probe accounting.

McCormick examples make the placement of exactness particularly visible.
The interaction graph identifies faces on which the separate bilinear gaps
vanish, and transverse near-optimal strata can force many boxes. The
aligned two-dimensional segment admits a two-leaf exact certificate, but the
aligned three-dimensional example still has a polynomial count. Thus graph
and tangent information supply useful sufficient tests and lower bounds;
they do not yet classify every aligned or curved case. Convergence order
alone omits this placement information.

For relaxed constraints, the obstruction comes from superoptimal infeasible
points retained by an inner relaxation tube. The $C^2$ objective and compatible original exact and relaxed constraint
data, LICQ, and the nonzero relaxed multiplier are substantive
hypotheses of the KKT argument. Exact original-feasibility reduction lies
outside its operation class because it can remove these witnesses directly.
Whether specific original-constraint contractors remove the obstruction on
curved active strata requires a separate analysis.

The decomposition comparison changes how local certificates are combined.
A spatial leaf assigns one global box. An affine-separator certificate can
combine local bag members without explicitly enumerating every global
combination. The fixed path separation in \cref{sec:decomposition} shows
that this can remove an exponential dependence on dimension even with the
same local factor relaxation and a unique nondegenerate minimizer. The
existence proof uses a known center and separator slopes; locating them and
checking intersecting local members have additional costs. It establishes a
change in proof representation, with explicit structural assumptions, rather
than a general algorithmic bound for every sparse problem.

For integer variables, the class number gives a corresponding convexity
obstruction. Every permitted leaf must contain an admissible class of the
required integer points. A large class number therefore persists under
convex cuts preserving the required integer set for the fixed projected
relaxation. Incumbent
removals instead enter the augmented count. Conversely, the example with
class number two and an exponential variable-branching tree demonstrates a
restriction of the disjunction class. The unrestricted semantic attainment
theorem identifies the minimum proof size; it supplies no efficient method
to construct the attaining convex pieces.

\subsection{Finding a proof and paying for its tests}\label{discussion:finding}

The one-dimensional minimizer theorem gives a case in which local relaxation
information suffices to find a near-minimal tree: its exact-gap rule is
4-competitive for every continuous objective, with the specified incumbent
convention. The information lower bounds have a different role. They use
smooth nonanalytic functions that agree near the observed minimizer, and so
apply to that stated information class. A real-analytic germ determines the
function on a connected domain; the same indistinguishability construction
cannot establish a lower bound there. In higher dimensions, all-coordinate splitting can lose a logarithmic
factor even on a separable quadratic. Repeated coordinate ambiguity forces
exponential dimension dependence for the specified node-local nonanalytic
information class with coordinatewise minimizer selection and no learning
from siblings. The sharp-coordinate theorem covers a restricted separable
class.

Safeguards impose a geometric constraint even if a rule knows where it
would prefer to cut. On kink families, an interior safety fraction can
prevent the desired cut for a long chain. The recentring rule attains the
minimum safe hitting time in this model; positive-tolerance pruning can
stop the chain sooner. The exact one-dimensional gap model also guarantees
that every relaxation minimizer at an invalid interval is strictly interior.
Production relaxations can attain their minimum on variable bounds. That
behavior, visible in the archived benchmark diagnostics, is a reason to
retain safeguards that the idealized kink model does not capture.

Propagation creates a similar distinction between the ideal test and its
implementation. The greatest hull-consistent fixed point defines the
strength of objective-only cutoff propagation for a specified expression
graph. Local exactness can yield a bounded ideal node count. The specified
expanded-quadratic schedule nevertheless needs polynomially many rounds in
$1/\epsilon$, while a particular lifted representation has a much smaller
round count. The persistence results require their derivative-cancellation
conditions and charge objective-only phases. Combining constraints,
auxiliary-variable tightening, or other contractors changes the test and
needs its own proof.

Probe tests and local decomposition checks are also part of finding or
verifying a certificate. A bound formed by the minimum of probe-child
bounds is justified by those child tests; a displayed one-node search does
not absorb their cost. Likewise, an unaligned local leaf may meet many
separator cells, each requiring a check. The accounting in this paper keeps
such work visible without claiming that a relaxation evaluation or a convex
program has uniform running time.

\subsection{What the statistical transitions distinguish}\label{discussion:statistical}

The sparse-regression linear-tree window proves that root inexactness can
coexist with a short certificate. Above the single-variable threshold, the
planted support is uniquely optimal and the C1 path lemma applies; the root
threshold occurs later. The converse compares relaxation bounds with the
planted-support value. It becomes a statement about global inexactness only
when that support is optimal. The low-total-signal certificate lower bound
uses a separate asymptotic regime, and the selected local-lift conclusions
retain their explicit sandwich and helper-column conventions.

Binary least squares separates three further events. The hard certificate
law concerns minimum proof size in the divergent sublinear-SNR regime. The sharp
linear-SNR transition concerns C1 at all single-wrong-fixing nodes, with
fixed margins; it is a threshold for a sufficient linear-tree criterion.
The logarithmic-SNR root transition concerns equality with the orthant
relaxation, so it says when the upper box constraints are inactive. It
implies neither an integral root optimum nor correct sign rounding.
Failure of C1 below its threshold likewise proves no tree lower bound on
its own. The finite NNLS coefficient estimate is useful because it controls
all node events simultaneously without assuming their independence.

The fixed-dimensional sparse-regression calculation also separates support
recovery from relaxation exactness. Under fixed positive per-entry noise and
ridge $\sqrt n$, the true support becomes optimal but the exactness
probability approaches a constant below one. Changing the noise
normalization changes this calculation. Such distinctions are necessary
when transferring a published recovery or exactness statement into a
branch-and-bound claim.

\subsection{What the archived computations support}\label{discussion:computations}

The archived computations in \cref{sec:experiments} illustrate the mechanisms
at their recorded finite sizes and tolerances. Under controlled settings
with a known optimal incumbent, best-first search, and heuristics switched
off, synthetic node counts follow the predicted positive-dimensional
exponents over the reported range. Those settings change several solver
components together. Default settings can change the effective relaxation
through propagation and reformulation, and numerical thresholds can mask
small-tolerance behavior. The fitted slopes describe floating-point runs;
they are not certified asymptotic exponents.

The benchmark comparisons provide no general gain from moving split points
closer to relaxation minimizers. Removing the clamp roughly doubled the
recorded node counts, while the recentering comparison had uncertainty
compatible with similar performance. The midpoint comparison showed a
modest node-count disadvantage in its separate recorded protocol. These
comparisons support the relevance of bound-attaining relaxation solutions
and safeguards, but do not rank branching rules across solver families.

Random-design transitions illustrate the difference between C1 and root
certification. Their sizes do not satisfy the easy sparse-regression
theorem's asymptotic conditions, and their observed locations do not estimate
its first-order constants. Some explicit hard-side constants are vacuous
at all recorded sizes. The computations therefore complement the proofs
with finite-size behavior and implementation effects. They establish no
general practical speedup, and node counts alone do not compare methods
whose node tests have different costs.

\subsection{Open questions}\label{discussion:open}

The following questions remain within reach of the certificate viewpoint.
\begin{enumerate}
\item Can a natural relaxation-minimizer coordinate rule be competitive in
  every fixed dimension under an exact-gap oracle? The present
  sharp-coordinate theorem leaves two arbitrary separable coordinates and
  the general nonseparable case unresolved. Dimension dependence must be
  assessed within the stated information class.
\item Can the graph and face bounds be completed into a characterization
  for termwise McCormick certificates? Matching integrals leave logarithmic
  losses, and the curved example lacks a matching upper bound. Aligned
  higher-dimensional examples rule out a classification based only on
  whether the tangent obstruction vanishes.
\item Which expression graphs admit local exactness, and which retain
  geometric lower bounds under cutoff propagation? The sufficient classes
  here leave a gap. Round complexity for general fair schedules, combined
  contractors, and other representations also remains unresolved.
\item Which original-constraint reductions remove the relaxed-constraint
  obstruction on curved active strata, and at what evaluation or round
  cost? The current tube theorem intentionally excludes those operations.
\item For convex quadratic projected objectives, how large can the minimum
  split-branching tree be relative to the class number? Semantic attainment
  by arbitrary convex pieces does not resolve that restriction.
\item What finite-size bounds describe the statistical transition windows?
  Between the proved divergent sublinear-SNR certificate law and the fixed-margin C1
  regime, the available statements do not give a complete size law. Failure
  of a sufficient test cannot fill that gap.
\item Can models that include relaxation minimizers on variable bounds
  explain the benchmark safeguard behavior and yield competitive guarantees
  with an explicit evaluation cost?
\end{enumerate}
